% I tre generi di leva, con il fulcro (triangolo), la forza motrice F_mot e la
% resistenza F_res, e un esempio nel corpo umano per ciascuno.
\begin{tikzpicture}[>=Stealth, font=\small]
  \tikzset{
    asta/.style={fill=richiamo!25, draw=richiamo, thick},
    fulcro/.style={fill=richiamo!70, draw=richiamo},
    motrice/.style={->, medica, very thick},
    resistenza/.style={->, attenzione, very thick},
  }
  % primo genere: fulcro fra le due forze
  \begin{scope}
    \filldraw[asta] (-1.6,0) rectangle (1.6,0.14);
    \filldraw[fulcro] (0.3,0) -- (0.05,-0.45) -- (0.55,-0.45) -- cycle;
    \draw[resistenza] (-1.4,0.14) -- (-1.4,-0.9) node[below] {$\vec F_{\text{res}}$};
    \draw[motrice] (1.4,0.14) -- (1.4,-0.9) node[below] {$\vec F_{\text{mot}}$};
    \node[align=center] at (0,1.9) {\textbf{primo genere}\\fulcro in mezzo};
    \node[align=center, text width=3.6cm] at (0,-2.35) {la testa sull'atlante: muscoli della nuca contro il peso del capo};
  \end{scope}
  % secondo genere: resistenza fra fulcro e forza motrice
  \begin{scope}[xshift=4.6cm]
    \filldraw[asta] (-1.6,0) rectangle (1.6,0.14);
    \filldraw[fulcro] (-1.5,0) -- (-1.75,-0.45) -- (-1.25,-0.45) -- cycle;
    \draw[resistenza] (-0.3,0.14) -- (-0.3,-0.9) node[below] {$\vec F_{\text{res}}$};
    \draw[motrice] (1.4,0.14) -- (1.4,1.1) node[right] {$\vec F_{\text{mot}}$};
    \node[align=center] at (0,1.9) {\textbf{secondo genere}\\resistenza in mezzo};
    \node[align=center, text width=3.6cm] at (0,-2.35) {il piede sulle punte: polpaccio contro il peso del corpo};
  \end{scope}
  % terzo genere: forza motrice fra fulcro e resistenza
  \begin{scope}[xshift=9.2cm]
    \filldraw[asta] (-1.6,0) rectangle (1.6,0.14);
    \filldraw[fulcro] (-1.5,0) -- (-1.75,-0.45) -- (-1.25,-0.45) -- cycle;
    \draw[motrice] (-0.8,0.14) -- (-0.8,1.1) node[right] {$\vec F_{\text{mot}}$};
    \draw[resistenza] (1.4,0.14) -- (1.4,-0.9) node[below] {$\vec F_{\text{res}}$};
    \node[align=center] at (0,1.9) {\textbf{terzo genere}\\forza motrice in mezzo};
    \node[align=center, text width=3.6cm] at (0,-2.35) {l'avambraccio: bicipite contro il peso in mano};
  \end{scope}
\end{tikzpicture}
